\documentclass[11pt]{article}
\usepackage[a4paper,margin=1in]{geometry}
\usepackage{amsmath,amssymb,mathtools,booktabs,enumitem}
\usepackage[dvipsnames]{xcolor}
\usepackage{hyperref,microtype,comment}
\hypersetup{colorlinks=true,linkcolor=MidnightBlue,urlcolor=MidnightBlue,
 pdftitle={Recitation 8: Midterm 2 Review},pdfauthor={Xuanxi Zhang}}
% Default: student version. Define \WITHSOLUTIONS for answers and pacing.
\newif\ifsolutions
\ifdefined\WITHSOLUTIONS\solutionstrue\else\solutionsfalse\fi
\newcommand{\rv}[1]{\widetilde{#1}}
\newcommand{\E}{\mathbb E}
\newcommand{\Pbb}{\mathbb P}
\newcommand{\Var}{\operatorname{Var}}
\newcommand{\Cov}{\operatorname{Cov}}
\newcommand{\Bias}{\operatorname{Bias}}
\newcommand{\MSE}{\operatorname{MSE}}
\newcommand{\se}{\operatorname{se}}
\ifsolutions
 \newenvironment{answer}{\par\medskip\noindent\textcolor{MidnightBlue}{\textbf{Solution.}}\ }{\par\medskip}
 \newcommand{\workingspace}[1]{}
\else
 \excludecomment{answer}
 \newcommand{\workingspace}[1]{\vspace{#1}}
\fi
\setlist[enumerate]{label=(\alph*),leftmargin=2em,itemsep=0.4em}
\setlist[itemize]{leftmargin=1.5em,itemsep=0.3em}
\title{\huge Recitation 8\\[0.35em]\normalsize Midterm 2 Review\\[0.35em]
 \large\ifsolutions Solutions and Teaching Notes\else Review and Practice Problems\fi}
\author{Xuanxi Zhang}
\date{October 30, 2026}
\begin{document}
\maketitle
\noindent\textit{Exam:} Tuesday, November 3, in class. Material after Midterm 1
through the bootstrap; no coding questions. Review textbook Chapters 7--9
alongside the lecture examples and assigned exercises.

\section*{Review: conditioning, prediction, and correlation}
\paragraph{Conditioning changes the population.}
For a discrete grouping variable $\rv{r}$,
\[
 \E[\rv{s}]=\sum_r\Pbb(\rv r=r)\E[\rv s\mid\rv r=r].
\]
Among respondents ($\rv q=1$), the weights instead become
\[
 \Pbb(\rv r=r\mid\rv q=1)
 =\frac{\Pbb(\rv q=1\mid\rv r=r)\Pbb(\rv r=r)}
 {\sum_u\Pbb(\rv q=1\mid\rv r=u)\Pbb(\rv r=u)}.
\]
If $\rv s$ and $\rv q$ are conditionally independent given $\rv r$,
the within-group mean is unchanged, but the overall observed mean can change.
\textbf{Check the weights before averaging.}

\paragraph{Best prediction under squared error.}
For finite second moments, the unrestricted minimum-MSE predictor is
\[
 g_*(y)=\E[\rv x\mid\rv y=y].
\]
The best affine predictor $\ell_*(y)=a+by$, when $\Var(\rv y)>0$, is
\[
 b=\frac{\Cov(\rv x,\rv y)}{\Var(\rv y)},\qquad
 a=\E[\rv x]-b\E[\rv y],\qquad
 \MSE(\ell_*)=\Var(\rv x)-\frac{\Cov(\rv x,\rv y)^2}{\Var(\rv y)}.
\]
The conditional mean need not be linear. Joint Gaussianity is one sufficient
condition for these two predictors to agree.

\paragraph{Covariance and correlation.}
\[
 \Cov(\rv x,\rv y)=\E[\rv x\rv y]-\E[\rv x]\E[\rv y],\qquad
 \rho_{\rv x,\rv y}=\frac{\Cov(\rv x,\rv y)}
 {\sqrt{\Var(\rv x)\Var(\rv y)}}.
\]
Correlation is defined when both variances are positive. Independence implies
zero covariance; the converse is false in general. Correlation alone does not
establish causation.

\paragraph{Average treatment effect.}
For potential outcomes $\rv w(1),\rv w(0)$,
$\mathrm{ATE}=\E[\rv w(1)-\rv w(0)]$.
Under consistency, conditional exchangeability given a covariate $\rv c$,
and positivity, adjust by averaging the within-stratum treatment contrasts
using the \emph{same target-population weights} for both groups.

\newpage
\section*{Review: sampling uncertainty and the bootstrap}
\paragraph{Observed data versus random estimators.}
For $X=\{x_1,\ldots,x_n\}$, the observed mean and sample variance are
\[
 m(X)=\frac1n\sum_i x_i,\qquad
 v(X)=\frac1{n-1}\sum_i(x_i-m(X))^2\quad(n\geq2).
\]
Before observing an i.i.d. sample with mean $\mu$ and variance $\sigma^2$,
the random sample mean $\rv m=n^{-1}\sum_i\rv x_i$ satisfies
$\E[\rv m]=\mu$ and $\se(\rv m)=\sigma/\sqrt n$.
The standard deviation describes individual observations; the standard error
describes an estimator across repeated samples.

\paragraph{Bias and variance measure different errors.}
For a fixed parameter $\theta$ and an estimator $\rv g$,
\[
 \Bias(\rv g)=\E[\rv g]-\theta,\qquad
 \se(\rv g)=\sqrt{\Var(\rv g)},\qquad
 \MSE(\rv g)=\Bias(\rv g)^2+\Var(\rv g).
\]
Increasing sample size does not automatically remove selection bias.
For deterministic weights $a_i$,
\[
 \Var\!\left(\sum_i a_i\rv m_i\right)
 =\sum_i a_i^2\Var(\rv m_i)
   +2\sum_{i<j}a_i a_j\Cov(\rv m_i,\rv m_j).
\]
Only discard the covariance terms when justified. If published estimates
reuse the same data, expand them in terms of the original independent estimates.

\paragraph{Confidence intervals.}
A $1-\alpha$ confidence procedure covers a fixed parameter in at least
$1-\alpha$ of repeated samples (approximately so for approximate procedures).
For an approximately unbiased, approximately Gaussian estimator, a common
95\% interval is
\[
 [g-1.96\se_{\mathrm{est}},\;g+1.96\se_{\mathrm{est}}].
\]
For the sample mean, $\se_{\mathrm{est}}=\sqrt{v(X)/n}$.
Higher confidence widens the interval; multiplying $n$ by four roughly halves
its width. The 95\% refers to the procedure's coverage, not a posterior
probability for the fixed parameter after seeing the interval.

\paragraph{Nonparametric bootstrap.}
Treat the empirical distribution of $X$ as a stand-in for the unknown population.
Draw each resample independently, taking $n$ observations \emph{with replacement}
from $X$, and recompute the same statistic. For $B$ bootstrap estimates
$g_1^*,\ldots,g_B^*$, write $\bar g^*=B^{-1}\sum_b g_b^*$. Then
\[
 \se_{\mathrm{boot}}=
 \sqrt{\frac1{B-1}\sum_{b=1}^B(g_b^*-\bar g^*)^2}.
\]
A percentile interval uses the empirical 2.5\% and 97.5\% quantiles of the
bootstrap estimates. Its accuracy depends on how well resampling approximates
the sampling distribution; it is not automatically an exact 95\% interval.
Increasing $B$ reduces simulation error, not uncertainty from a small original sample.

\newpage
\section*{Practice problems}
\noindent Three core problems from the official Midterm 2 practice set;
Problem 4 is a backup. Work independently before discussion.

\subsection*{Problem 1: Hotel questionnaire (official Problem 1)}
\noindent\textit{Independent work: 7 minutes.}
A hotel asks customers for a discrete satisfaction score between 1 (worst)
and 5 (best). Twenty percent of customers are retired. Retired customers
have mean score 2 and respond with probability 0.9. Other customers have
mean score 4 and respond with probability 0.25.

Let $\rv r$ indicate retirement, $\rv q$ indicate responding, and $\rv s$
denote the score. Assume $\rv q$ and $\rv s$ are conditionally independent
given $\rv r$.
\begin{enumerate}
 \item Find $\E[\rv s]$ for the full customer population.
 \item Find the mean observed score $\E[\rv s\mid\rv q=1]$.
\end{enumerate}
\noindent\textit{Discussion:} Which group is overrepresented among respondents?
Would collecting more responses in the same way fix the discrepancy?
\workingspace{9cm}
\begin{answer}
\textbf{(a)} By iterated expectation,
\[
 \boxed{\E[\rv s]=0.2(2)+0.8(4)=3.6.}
\]
\textbf{(b)} The response probability is
$\Pbb(\rv q=1)=0.2(0.9)+0.8(0.25)=0.38$.
Bayes' rule gives the composition of the responding group:
\[
 \Pbb(\rv r=1\mid\rv q=1)=\frac{0.18}{0.38}=\frac9{19},\qquad
 \Pbb(\rv r=0\mid\rv q=1)=\frac{10}{19}.
\]
Conditional independence implies
$\E[\rv s\mid\rv r=r,\rv q=1]=\E[\rv s\mid\rv r=r]$.
Consequently,
\[
 \boxed{\E[\rv s\mid\rv q=1]
 =\frac9{19}(2)+\frac{10}{19}(4)=\frac{58}{19}\approx3.053.}
\]
Retired customers make up about 47.4\% of respondents but only 20\% of
customers. Their lower mean score pulls the respondent mean downward by
$3.6-58/19=52/95\approx0.547$.

More independent responses from the same respondent distribution reduce
standard error but leave this population discrepancy. If the population
group shares are known, weighting the two respondent group means by 0.2
and 0.8 targets the population mean under the stated conditional-independence assumption.
\end{answer}

\newpage
\subsection*{Problem 2: Noisy measurement (official Problem 6)}
\noindent\textit{Independent work: 8 minutes.}
We want to measure a quantity $\rv x$ with mean zero and variance one.
We observe $\rv y=\rv x+\rv z$, where the noise $\rv z$ has mean zero,
variance $\sigma^2$, and is independent of $\rv x$.
\begin{enumerate}
 \item Find the best linear (affine) estimator of $\rv x$ given $\rv y=y$.
 \item Compute its mean squared error.
 \item What happens to the estimator and its error as $\sigma\to0$?
 Explain why this makes sense.
 \item What happens as $\sigma\to\infty$? Explain why this makes sense.
\end{enumerate}
\noindent\textit{Discussion:} Do the assumptions guarantee that the best linear
predictor equals the unrestricted conditional-mean predictor?
\workingspace{10cm}
\begin{answer}
\textbf{(a)} Independence and zero means give
\[
 \E[\rv y]=0,\qquad \Var(\rv y)=1+\sigma^2,\qquad
 \Cov(\rv x,\rv y)=\Var(\rv x)+\Cov(\rv x,\rv z)=1.
\]
The intercept is zero and the slope is $1/(1+\sigma^2)$, so
\[
 \boxed{\ell_*(y)=\frac{y}{1+\sigma^2}.}
\]
\textbf{(b)} Set $b=1/(1+\sigma^2)$. Since
$\rv x-b\rv y=(1-b)\rv x-b\rv z$, independence gives
\[
 \MSE(\ell_*)=(1-b)^2+b^2\sigma^2
 =\frac{\sigma^4+\sigma^2}{(1+\sigma^2)^2}
 =\boxed{\frac{\sigma^2}{1+\sigma^2}}.
\]
Equivalently, minimize the quadratic
$\E[(\rv x-b\rv y)^2]=1-2b+b^2(1+\sigma^2)$ directly.

\textbf{(c)} As $\sigma\to0$, $\ell_*(y)\to y$ and the MSE tends to zero.
With no noise, the measurement equals the quantity of interest.

\textbf{(d)} As $\sigma\to\infty$, the slope tends to zero and the MSE tends
to one. The linear estimator puts vanishing weight on the increasingly noisy
measurement and approaches the prior mean, zero. Its random output also
has variance $\Var(\ell_*(\rv y))=1/(1+\sigma^2)\to0$.

\textbf{Discussion.} No Gaussian assumption was given, so this is optimal
within the affine class only. The unrestricted optimum is
$\E[\rv x\mid\rv y=y]$, which can be nonlinear.
\end{answer}

\newpage
\subsection*{Problem 3: Herding (official Problem 10)}
\noindent\textit{Independent work: 10 minutes.}
Three independent polls produce unbiased estimators
$\rv m_1,\rv m_2,\rv m_3$ of a fixed support proportion $\theta$.
Each has standard error $s$. After seeing earlier published polls,
pollsters modify their estimates by averaging:
\[
 \rv h_2=\frac{\rv m_1+\rv m_2}{2},\qquad
 \rv h_3=\frac12\left(\frac{\rv m_1+\rv h_2}{2}+\rv m_3\right).
\]
\begin{enumerate}
 \item Are $\rv h_2$ and $\rv h_3$ unbiased? Does herding introduce bias here?
 \item Find their standard errors. Does herding improve each individual
 estimator under this model?
 \item Compare the standard errors of
 \[
 \rv m_{\mathrm{all}}=\frac{\rv m_1+\rv m_2+\rv m_3}{3},\qquad
 \rv h_{\mathrm{all}}=\frac{\rv m_1+\rv h_2+\rv h_3}{3}.
 \]
 Why is herding problematic when we combine the published polls?
\end{enumerate}
\noindent\textit{Hint:} First express each estimator in terms of the original
independent polls $\rv m_1,\rv m_2,\rv m_3$.
\workingspace{8cm}
\begin{answer}
\textbf{(a)} Expand
$\rv h_3=\tfrac38\rv m_1+\tfrac18\rv m_2+\tfrac12\rv m_3$.
The weights of both $\rv h_2$ and $\rv h_3$ sum to one, so
$\E[\rv h_2]=\E[\rv h_3]=\theta$: both remain unbiased.

\textbf{(b)} Apply independence to the \emph{original} polls:
\[
 \Var(\rv h_2)=\frac{s^2}{2},\qquad
 \Var(\rv h_3)=\left(\frac9{64}+\frac1{64}+\frac14\right)s^2
 =\frac{13}{32}s^2.
\]
Thus $\se(\rv h_2)=s/\sqrt2\approx0.707s$ and
$\se(\rv h_3)=\sqrt{13/32}\,s\approx0.637s$.
Each has smaller variance and MSE than its corresponding original poll,
because it borrows information from other polls.

\textbf{(c)} The unmodified average has variance $s^2/3$.
For the published, herded polls,
\[
 \rv h_{\mathrm{all}}=
 \frac58\rv m_1+\frac5{24}\rv m_2+\frac16\rv m_3,
 \qquad
 \Var(\rv h_{\mathrm{all}})
 =\frac{225+25+16}{576}s^2=\frac{133}{288}s^2.
\]
Therefore
\[
 \boxed{\se(\rv m_{\mathrm{all}})=\frac{s}{\sqrt3}\approx0.577s,
 \quad \se(\rv h_{\mathrm{all}})=\sqrt{\frac{133}{288}}s\approx0.680s.}
\]
Both averages are unbiased, but herding increases the aggregate variance by
a factor $133/96\approx1.385$. It overweights the first poll and creates
dependence among published estimates; treating them as independent would be wrong.
\end{answer}

\newpage
\subsection*{Problem 4: A tiny bootstrap (backup)}
\noindent\textit{Independent work: 5 minutes. Adapted from archived Recitation
Worksheet 8, Problem 1.}
The observed dataset is $X=\{2,4\}$. For a sample of size two, define the
median as the average of the two ordered observations.
\begin{enumerate}
 \item List all ordered bootstrap resamples of size two, drawn with replacement.
 Give each resample's probability.
 \item Compute the median of each resample and the exact bootstrap mean.
 \item Compute the exact bootstrap variance and standard error of the median.
 Use the probability distribution from (a), rather than a Monte Carlo estimate.
\end{enumerate}
\noindent\textit{Discussion:} Why does enumerating all possible resamples use
probability weights rather than the $B-1$ formula for simulated replicates?
\workingspace{10cm}
\begin{answer}
\textbf{(a)--(b)} Each draw independently selects 2 or 4 with probability $1/2$.
\begin{center}
\begin{tabular}{@{}ccc@{}}
\toprule Resample & Probability & Median\\\midrule
$(2,2)$ & $1/4$ & 2\\
$(2,4)$ & $1/4$ & 3\\
$(4,2)$ & $1/4$ & 3\\
$(4,4)$ & $1/4$ & 4\\\bottomrule
\end{tabular}
\end{center}
The bootstrap mean is $(2+3+3+4)/4=3$, equal to the observed median.

\textbf{(c)} Conditional on the fixed original data, the exact bootstrap
variance and SE are
\[
 \boxed{\Var_*(\rv g^*)=\frac{(2-3)^2+(3-3)^2+(3-3)^2+(4-3)^2}{4}
 =\frac12,\qquad \se_*=\frac1{\sqrt2}.}
\]
Here the four entries enumerate the entire resampling distribution with
known probabilities; they are not four independently simulated replicates.
If we simulate $B$ replicates instead, the sample variance with denominator
$B-1$ estimates this bootstrap-distribution variance.
This exact bootstrap calculation is still only a plug-in estimate of the
unknown population sampling variance. Two original observations give little
information about that population.
\end{answer}

\ifsolutions
\newpage
\section*{Instructor notes: 75-minute plan}
\begin{center}
\begin{tabular}{@{}p{0.16\textwidth}p{0.77\textwidth}@{}}
\toprule Time & Activity\\\midrule
0--10 min & Conditioning, prediction, correlation, and ATE; ask students to
 distinguish a subgroup mean from the population mean.\\
10--25 min & Bias versus SE; covariance of averages; confidence-interval
 interpretation; resampling with replacement.\\
25--39 min & Problem 1: 7 minutes independent work, 7 minutes discussion.\\
39--54 min & Problem 2: 8 minutes independent work, 7 minutes discussion.\\
54--70 min & Problem 3: 10 minutes independent work, 6 minutes discussion.\\
70--75 min & Exit checks and questions. For a fast class, save 3 minutes
 during earlier discussions: use 67--72 for Problem 4 and 72--75 to discuss it.\\
\bottomrule
\end{tabular}
\end{center}

\paragraph{Why these problems?}
The October 29 review deck contains official practice Problems 3
(weight-loss supplement), 4 (three random variables), 11 (confidence-interval
plots), 12 (cholesterol), and 13 (modified Gaussian bootstrap).
Problems 1, 6, and 10 do not appear in that deck and form the core here.
This is a comparison with the posted slides, not a claim that every displayed
problem was actually completed in the lecture. Problem 13 appears only as
a prompt on the last slide, so its classroom completion is especially uncertain.

\paragraph{Sources and scope.}
\begin{itemize}
 \item \href{https://brightspace.nyu.edu/d2l/le/lessons/613220/units/13329394}
 {Brightspace: Schedule and Course Material}, October 29 row:
 \emph{Practice Problems for Midterm 2}, Problems 1, 6, and 10;
 corresponding official solutions; \emph{Review}, 33 PDF pages.
 Local reference copies are stored in \texttt{r8/sources/}.
 \item Backup: \emph{Recitation Worksheet 8}, Problem 1, in the user's 2025
 archive (the worksheet itself is dated March 21, 2025).
 Wording clarifies exact enumeration versus Monte Carlo variance estimation.
 \item The October 27 exam announcement specifies November 3 and material
 after Midterm 1 through the bootstrap, without coding questions.
\end{itemize}

\paragraph{Corrections to the official Herding solution.}
Equation (90) should read $\theta_{\mathrm{pop}}$, not
$\theta_{\mathrm{pop}}/2$. Equation (101) should have coefficients
$(15,5,4)/24$ on $(\rv m_1,\rv m_2,\rv m_3)$, respectively.
In equation (102), the third variance is $\Var(\rv m_3)$.
In equation (104), the first squared numerator is $15^2$, not $13^2$.
The official final variance $133s^2/288$ is correct.

\paragraph{Exit checks.}
Can students explain why more respondents do not remove response bias,
why the best linear predictor need not be the conditional mean, and why
averaging dependent published polls can lose information?
\fi
\end{document}
